\newcounter{counterboxspe}
\setcounter{counterboxspe}{0}

\newcounter{counterboxspetosave}
\setcounter{counterboxspetosave}{0}
\stepcounter{counterboxspetosave}
\newcounter{Empreinte}
\setcounter{Empreinte}{\thecounterboxspetosave{}}
\stepcounter{counterboxspetosave}
\newcounter{GT}
\setcounter{GT}{\thecounterboxspetosave{}}
\stepcounter{counterboxspetosave}
\newcounter{Collaborative}
\setcounter{Collaborative}{\thecounterboxspetosave{}}
\stepcounter{counterboxspetosave}
\newcounter{Conferences}
\setcounter{Conferences}{\thecounterboxspetosave{}}
\stepcounter{counterboxspetosave}
\newcounter{Publication}
\setcounter{Publication}{\thecounterboxspetosave{}}



\addchapnonumber{A\ \ Témoignage personnel}\label{annex_temoignage}

%\chapter{Témoignage personnel}\label{annex_temoignage}
%\chaptermark{\color{curcolor} Témoignage personnel}



\poetry{8}{

Si je dois me tenir, alors je prendrai la parole;

Si je dois me taire, alors je me lèverai.

~

Dire tout haut ce qui se murmure,

Faire sans bruit ce qui fait fracas.

}


\newpage






%%%%%%%%%%%%%%%%%%%%%%%%%%%%%%%%%%%%%%%%%%%%%%
{\noindent\Large\textbf{Préambule}}
\\

La recherche est un incroyable espace de liberté et de créativité, mais c'est aussi le paradis du questionnement, de la remise en question, de la critique. Il me semble donc intéressant d'avoir une approche critique sur le système global, unique et codifié qu'est celui de la recherche. Parce qu'une thèse est une formation, alors j'essaierai de tirer quelques leçons de l'expérience que j'en ai eue. Ici je me livre, à titre personnel, à travers des éléments de réflexion qui m'ont accompagné au cours de cette thèse, car je considère que nous avons besoin d'initier une telle réflexion sur le système de la recherche à l'échelle individuelle et collective, sans plus attendre.
\\


%%%%%%%%%%%%%%%%%%%%%%%%%%%%%%%%%%%%%%%%%%%%%%
{\noindent\Large\textbf{Débattre}}
\\

La recherche me semble rongée par l'immatériel, qu'il soit économique, informationnel, ou encore administratif. En particulier, il me paraît essentiel d'accepter, et donc de permettre de soulever la problématique, que le cadre de la recherche s'étend au-delà de l'activité de recherche en tant que telle, la sortant de son cadre idéalisé, rationnel et objectif. Car il n'y a pas de recherche sans moyens mis en oeuvre, et il n'y a pas non plus de savoir collectif ou encore d'application sans partage. La recherche fait partie intégrante de la société, et elle en est fondamentalement intriquée. Elle s'abreuve et alimente tant la sphère décisionnelle, à travers la politique, que la sphère applicative, à travers le lien avec le secteur privé. Enfin, elle nourrit une immense sphère didactique qui vise à la construction de l'édifice de savoirs ainsi qu'à les enseigner ou encore les communiquer au grand public. Bien qu'on s'éloigne ici de la méthode scientifique, cela fait fondamentalement partie du cadre de la recherche, et mérite donc qu'on en débatte. Et il y a beaucoup à débattre.


Tout d'abord, face aux crises sociales et environnementales, notamment la question du changement climatique et de l'effondrement du vivant, notre activité de recherche ne devrait-elle pas être repensée à la hauteur des enjeux ? Partager sa recherche, c'est communiquer un message, mais ici, quel est donc notre message, en particulier pour nous, scientifiques sur le climat ? Où allons-nous, que voulons-nous ? Sommes-nous happés par la vitesse de la recherche et ses contraintes pour vraiment se poser cette question ? Je pense notamment à l'empreinte environnementale de la recherche, dont l'empreinte carbone fait partie (voir Box~\hyperref[box:Empreinte]{\uppercase\expandafter{A\relax}.\theEmpreinte{}}). Je m'interroge par exemple en sciences du climat sur le fait de porter des projets de recherche qui nécessitent d'aller sur le terrain à l'autre bout du monde, et qui n'ont donc rien de durable. En particulier, si personne plus proche que soi du lieu d'étude n'est en mesure de réaliser une telle mission sur le terrain, c'est qu'il y a probablement un souci d'inclusion dans la recherche, n'est-il pas ? Mais de manière plus globale, je me questionne sur l'utilité de savoir pour ne pas adapter notre manière de faire en fonction de ce que l'on sait, si ce n'est nourrir, un jour très probablement, des regrets. Pour moi, la recherche en sciences du climat a ici un rôle majeur d'exemplarité à jouer, mais cela nécessite une profonde remise en question du système de la recherche. Radicale.

Pour le reste, il y a beaucoup à dire, je crois, et ce ne sont pas les échanges entre jeunes chercheu.r.se.s notamment qui feront taire ces remarques (voir Box~\hyperref[box:GT]{\uppercase\expandafter{A\relax}.\theGT{}}). Car la recherche, c'est une hiérarchie, invisible vue de l'extérieur, mais pénétrante, et parfois pesante, vu de l'intérieur. Des statuts définis à l'échelle globale que l'on pourrait très certainement interroger, et qui ne sont pas sans référence au caractère compétitif de la recherche. Loin de ce qu'elle pourrait laisser transparaître, la recherche n'est pas une unité soudée, c'est une somme d'individualités en compétition à l'échelle individuelle. Le paradoxe réside dans le fait que l'on ait besoin des autres pour mener à bien sa recherche tout en étant évalué individuellement. Cela crée un lien de dépendance intéressé et non collaboratif, à l'image du caractère séquentiel de la recherche, où l'on ne travaille pas ensemble, mais les uns à la suite des autres (voir Box~\hyperref[box:Collaborative]{\uppercase\expandafter{A\relax}.\theCollaborative{}}).


\begin{spacing}{1.5}
~
\end{spacing}


\boxspenoref{\thecounterboxspe{}}{L'empreinte de ma recherche}{
\setlength\parindent{12pt}
Il y a des traces que l'on veut faire ou laisser, et des traces que l'on aimerait rendre discrètes ou effacer. L'empreinte environnementale de la recherche en est justement une que l'on aimerait rendre discrète, non pas en la cachant, mais en essayant de la limiter au maximum. Bien que je ne sois pas un fervant défenseur du partage des empreintes carbones individuelles\footnote{Je considère qu'une empreinte carbone individuelle est capitale et informative pour soi, mais avec un faible intérêt pour les autres, contrairement à des empreintes plus globales, comme à l'échelle d'un laboratoire. Le partage d'une empreinte carbone à l'échelle individuelle amène une forme de transparence, mais je n'aime pas cette possibilité de pointer du doigt des individus.}, je souhaitais tout de même partager l'empreinte carbone relative à mon activité de recherche au cours de ces trois années de thèse, car elle est illustrative de mon engagement d'une part, mais aussi de la part conséquente de mon activité principale, à première vue incompressible, la modélisation. L'empreinte carbone de cette thèse s'élève ainsi à environ 4,3~tCO\textsubscript{2eq}\footnote{Elle a été réalisée à partir de chiffres communiqués en novembre 2022 par Olivier Aumont, et notamment la considération qu'une heure CPU équivaut à 1,1~gCO\textsubscript{2eq}.}, dont près de 35~\% est associée aux calculs numériques réalisés sur des super-calculateurs dans le cadre de la réalisation de simulations numériques.

\begin{spacing}{0.7}
~
\end{spacing}

\centerline{\begin{minipage}{0.85\textwidth}
\centering\includegraphics[width=0.5\textwidth]{others/Empreinte.png}
\end{minipage}}

\begin{spacing}{0.7}
~
\end{spacing}

Je souhaitais aussi revenir sur la part relativement faible des transports dans l'empreinte carbone de ma recherche. C'est sur ce point que mon engagement individuel s'est porté au cours de ma thèse en décidant de ne pas prendre l'avion ni pour aller en conférences, ni pour aller sur du terrain. Ainsi, je me déplace uniquement en train, en bus, et j'ai décidé de le mentionner maintenant, que ce soit dans mes présentations ou mes posters, avec un petit message à la clé :

\begin{spacing}{0.7}
~
\end{spacing}

\centerline{\begin{minipage}{0.85\textwidth}
\centering\includegraphics[width=0.7\textwidth]{others/Carbon_footprint.png}
\end{minipage}}

\begin{spacing}{0.7}
~
\end{spacing}

D'ailleurs, je crois que mon histoire peut être portée comme un message d'espoir, car j'ai passé dix ans de mon enfance à rêver et à vivre pour devenir pilote de ligne, avant que, pour des raisons environnementales, économiques et sociales, je ne décide d'abandonner ce rêve qui m'obnubilait. Alors si je suis capable de ne plus prendre l'avion, tout le monde en est capable, j'en suis persuadé. Mais il faut aussi que d'un point administratif cela suive, et vite. Réserver un train de nuit pour aller à Vienne est une épreuve, il faut le vouloir, il faut se battre. De même que prendre un 'Pass Interrail' pour se rendre en conférence en train à prix avantageux demande non seulement d'avancer la somme pendant de nombreux mois en attendant le remboursement, mais aussi une pile de justificatifs hallucinante, là où des billets d'avion ne demanderait quasi aucun effort. Sortir du cadre défini demande de l'énergie, et je souhaiterais que n'importe qui, plus tard, puisse prendre un tel passe sans difficulté.

La considération que j'apporte à mon empreinte carbone m'a notamment amené à refuser la proposition de participer à une campagne océanographique avec un départ d'Afrique du Sud, car je trouve cela
\label{box:Empreinte}
}


\boxspecontinuednoref{\thecounterboxspe{}}{L'empreinte de ma recherche}{
\setlength\parindent{12pt}
\noindent aberrant quand des Sud-africains pourraient mener une telle campagne plutôt que des Hollandais. De plus, à l'heure actuelle, je préfère ne pas m'engager pour des missions sur le terrain qui ne correspondent pas à mes convictions, celles d'un travail sur le terrain à faible empreinte environnementale, qui, dans le domaine des géosciences, se traduit généralement par une expérience au plus près de la nature, une aventure. Enfin, il est important de noter que l'empreinte carbone ne représente qu'un élément de l'empreinte environnementale totale, un fait que je garde également à l'esprit.
}


\begin{spacing}{1.5}
~
\end{spacing}






\boxspenoref{\thecounterboxspe{}}{Mon investissement dans des groupes de travail}{
\setlength\parindent{12pt}
J'ai très rapidement su que je souhaitais m'engager au cours de cette thèse, mais j'ai longuement hésité sur la manière dont je pouvais m'engager, et je me questionne toujours beaucoup à ce sujet. Ma réflexion n'est pas encore aboutie, je le sais, et dans tous les cas, ma pensée restera évolutive. J'ai initialement pensé rejoindre le collectif qui oeuvre pour sensibiliser et repenser notre empreinte carbone au laboratoire, 'LMD-Climaction', mais j'avais en réalité envie de me lancer dans une démarche plus pérenne, détachée de mon laboratoire. Alors, je me suis investi dans 'Labos 1point5' (\url{https://labos1point5.org}; dernier accès: Septembre 2023), un collectif de membres du monde académique, de toutes disciplines et sur tout le territoire, partageant un objectif commun : mieux comprendre et réduire l'impact des activités de recherche scientifique sur l'environnement, en particulier sur le climat. J'ai fait partie d'un groupe de travail sur les méthodes et pratiques de publication, et je fais maintenant partie d'un groupe de travail de jeunes chercheu.r.se.s. Je fais aussi partie de 'Scientifiques en rébellion' (\url{https://scientifiquesenrebellion.fr}; dernier accès: Septembre 2023), un collectif de scientifiques de toutes disciplines, tous statuts et tous secteurs qui se mobilisent contre l'inaction face au dérèglement climatique et à l'effondrement de la biodiversité. Je suis néanmoins passif à l'heure actuelle au sein de ce groupe, car j'ai besoin de réfléchir encore sur le mode d'action qui me correspond le mieux pour m'investir. L'élan de ces collectifs est merveilleux. Il nous permet de prendre consciences que l'on n'est pas seul, et qu'ensemble, il y a de belles choses à aller chercher.
\label{box:GT}
}

\begin{spacing}{1.5}
~
\end{spacing}

La compétition n'est pas forcément un problème en tant que tel, mais elle le devient lorsque le cadre disciplinaire est incohérent, et fondamentalement injuste. C'est, je crois, le cas de la recherche, dont le système d'évaluation est grotesque, à l'image du manque de nuance que peut avoir la recherche en certains points aujourd'hui. Le système d'évaluation de la recherche est aussi absurde que l'évaluation d'une performance sportive à l'applaudimètre. Ne pouvons-nous pas admettre, simplement reconnaître, que le travail des chercheu.r.se.s ne peut être comparé de manière quantitative ? Car tant de choses ne dépendent pas de soi, tant de choses ne dépendent pas de son investissement, de son effort. Considérer que l'on peut évaluer les chercheu.r.se.s de manière quantitative, c'est mépriser le travail de la recherche, et réifier les chercheu.r.se.s, de la même manière que nous avons aussi tendance à le faire avec le vivant. Or, je crois que le qualitatif peut tout autant être raisonné. Il est néanmoins plus nuancé, à l'écoute de la diversité, et il demande, en revanche, fondamentalement plus d'effort pour le développer. Le grotesque dans les indicateurs qui sont utilisés pour évaluer la recherche est que, d'une part, ils ne considèrent que le produit fini, la publication, et non le chemin parcouru pour y parvenir, et d'autre part, ils ne justifient en rien la qualité de la recherche produite, ni son apport pour la communauté. En fait, ces indicateurs retranscrivent uniquement la perception de la recherche à travers le prisme du système de l'information, qui sort du rationnel. Est-ce que la finalité de la recherche, c'est de publier beaucoup et d'être cité ? Et est-ce qu'être cité, c'est faire de la bonne recherche ? Ces indicateurs ont le mérite de détourner l'objetif de la recherche, car une fois défini, un indicateur, utilisé comme un outil d'évaluation de la performance d'individus à des fins sélectives, que ce soit pour avoir un poste permanent, ou encore avoir des financements, est voué à être considéré en premier lieu par les concerné.e.s, ici les chercheu.r.se.s. Avec de tels indicateurs, nul doute qu'un fossé puisse se creuser entre la qualité de la recherche et sa visibilité, car la mécanique du système de l'information exaltera toujours les sujets qui font du bruit.


\begin{spacing}{1.5}
~
\end{spacing}

\boxspenoref{\thecounterboxspe{}}{D'une somme d'individualité à un collectif}{
\setlength\parindent{12pt}
Au cours de cette thèse, et notamment au travers de mon premier travail d'intercomparaison des ESMs, je me suis beaucoup interrogé sur le caractère nominatif de la recherche alors qu'il se cache en réalité un effort collectif. En particulier, je n'aime pas du tout le fait que l'on cite les papiers, les rapport, les livres par le nom de l'auteur.e principal.e. Je trouve que cela exalte les individualités au détriment du collectif, car malgré le côté séquentiel plus que collaboratif de la recherche, nos travaux sont généralement le fruit d'un travail collectif. Pour moi, ce n'est pas une simple liste de co-auteurs qui matérialise un tel travail, c'est le partage sous la même banière du produit fini. Alors j'ai essayé, sans succès, de signer mon premier papier avec un nom de groupe en y glissant un petit message derrière :

\begin{spacing}{0.7}
~
\end{spacing}

\centerline{\begin{minipage}{0.85\textwidth}
\setlength\parindent{12pt}
\textit{The Marine Biogeochemistry Modelling Network (TMBMN): A collective, open and inclusive signature for marine biogeochemistry modellers to work towards thinking and building sustainable, well-reasoned and responsible research.
\newline We stand for collaboration, criticism, and commitment rather than competitiveness, assessment and individualism in research. Let us erase personal interests in favour of disinterested, but ethical research.
\newline Let us foster cooperations and synergies even outside current funding constraints.
\newline Let us develop an approach of transparency, accessibility and sharing.
\newline Let us reply to our relative ignorance with humility by giving life to our curiosity.}
\end{minipage}}

\begin{spacing}{0.7}
~
\end{spacing}

Signer avec ce nom de groupe aurait probablement pu être possible, mais cela aurait nécessité un investissement conséquent pour lui donner une voix, qu'il soit réutilisé par la suite dans le cadre d'études similaires. Je n'ai pas pris l'initiative de m'investir dans cette entreprise étant donné le peu de temps qu'il me restait pour finir ma thèse, mais je suis satisfait d'avoir un peu perturbé certaines personnes avec cette initiative, et que cela soit remonté jusqu'aux éditeurs en chef de \textit{Biogeosciences}.

C'est pour ces raisons aussi que j'ai pris l'initiative de modifier dans mon manuscrit, dans les chapitres qui ne correspondent pas à des papiers, l'ensemble des références aux rapports du GCB, de l'IPCC et de l'IPBES, de telle sorte que ce ne soient pas les auteur.e.s princip.aux.ales qui soient cité.e.s mais bien le nom des groupes respectifs.

\label{box:Collaborative}
}

\begin{spacing}{1.5}
~
\end{spacing}

 
Et en parlant de visibilité, il y en a une sur laquelle je me questionne, beaucoup, celle de l'utilité des réseaux sociaux dans le partage de connaissances (voir Box~\hyperref[box:Conferences]{\uppercase\expandafter{A\relax}.\theConferences{}}). Est-ce véritablement la bonne manière de communiquer ? Est-ce que cela n'est pas dévalorisant de partager au même endroit, avec la même valeur et la même visibilité, le fruit d'un consensus scientifique sur plusieurs décennies et une remarque climatosceptique que n'importe qui pourrait sortir en une fraction de seconde, que ce soit pour rire, ou de manière sérieuse ? N'est-ce pas à l'image d'une recherche qui se cherche dans sa manière de communiquer au grand public, face notamment à un système de l'information en transformation ? Mais à quoi bon rentrer dans le jeu de l'information qui se consumme, bien souvent superficielle et sans nuance, qui inonde, qui nous rend transparent à cette même information, qui nous transperce sans qu'on la digère.

\begin{spacing}{1.0}
~
\end{spacing}

\boxspenoref{\thecounterboxspe{}}{L'organisation de cycles de conférences}{
\setlength\parindent{12pt}
L'initiation d'échanges autour du climat et de l'Anthropocène, mais aussi autour de notre de manière de partager et d'agir en tant que scientifiques du climat, me semble clé dans mon processus de réflexion, et c'est aussi dans cette direction que j'aime communiquer. Avec Antoine Ehret et Rémi Gaillard, en partenariat avec des jeunes de l'association 'PC Durable' de l'ESPCI Paris - PSL, nous avons ainsi organisé un cycle de conférences à l'ESPCI Paris - PSL, et en partenariat avec l'IPSL, intitulé : "Climat \& Anthropocène", avec les interventions de Jean Jouzel, Corinne Le Quéré, Laurent Bopp, Juliette Mignot, Amaëlle Landais, Céline Guivarch, Jean-Baptiste Fressoz et Fabrice Flipo. En particulier, outre la volonté de partager aux étudiant.e.s de l'école, Antoine, Rémi et moi, souhaitions envoyer un message fort à l'administration de l'école par laquelle nous sommes passés pour une refonte de l'enseignement en prenant en considération les enjeux environnementaux et sociaux contemporains. Car nos échanges avec l'administration étaient restés jusqu'alors vains, ou presque.


\begin{spacing}{0.7}
~
\end{spacing}

\centerline{\begin{minipage}{0.85\textwidth}
\centering\includegraphics[width=0.9\textwidth]{others/Conférences.png}
\end{minipage}}

\begin{spacing}{0.7}
~
\end{spacing}

En fin de thèse, je me suis lancé dans l'organisation d'un second cycle de conférences au département de Géosciences de l'ENS - PSL, et en partenariat avec l'IPSL, intitulé : "Partager \& Agir", avec les interventions à venir de Anna Vayness, Laurent Bopp, Florence Habets, Camille Etienne, Sétphane Foucard, Fabienne Barataud, Odin Marc, Christophe Cassou, Audrey Cerdan, Paloma Moritz et Thomas Wagner. Une manière pour moi de proposer que l'on s'accorde du temps au sein de la communauté de climatologues pour échanger sur nos manières de partager et d'agir autour de notre sujet de recherche, ô combien médiatique et politique de nos jours.

\label{box:Conferences}
}


\begin{spacing}{1.5}
~
\end{spacing}

Et que dire aussi du financement de la recherche qui tend à contraindre de plus en plus sa temporalité, à travers des financements sur projet et des contrats courts ? Car nul ne peut savoir où mènera exactement la recherche, elle est faite d'incertitudes, de surprises, d'obstacles. La recherche est peut-être autant pilotée par la prise de risque que par celle d'en avoir une vision. La contraindre, c'est probablement l'accélérer au détriment de sa qualité, en la fractionnant, et lui ôtant certaines nuances, que le temps, seul, permet de soulever. C'est aussi mépriser l'inertie inhérente au processus de publication scientifique. Imposer une temporalité à la recherche, c'est pousser à répondre à des questions plutôt que de répondre à des problèmes, ouvrant beaucoup de portes sans véritablement les fermer finalement, car c'est, pour moi, en prolongeant sa recherche que l'on se rend compte de la limite de ses résultats. Plus on en fait, plus on a de chance de nuancer ses résultats, plus on doute, mais plus on comprend de manière générale. J'ai aussi un doute sur le fait que le système actuel de la recherche favorise l'émergence de nouvelles idées, car faire émerger quelque chose de nouveau demande un effort considérable pour sortir du cadre classique, ce que la temporalité de la recherche actuelle me paraît freiner. Au contraire, et c'est paradoxal, elle semble l'accélérer lorsqu'elle est pieds et poings liée avec le secteur privé. Ce même secteur qui de manière absurde porte le fruit de notre recherche, nos publications scientifiques (voir Box~\hyperref[box:Publication]{\uppercase\expandafter{A\relax}.\thePublication{}}). Comment peut-on accepter que du profit soit tiré des publications scientifiques ? Comment peut-on accepter que les chercheu.r.se.s paient à la soumission plus que le prix coûtant des corrections et de la mise en page ? N'est-il pas temps de reconsidérer globalement, et en profondeur, le système de publication~?

Je crois que la recherche mérite d'être repensée, redéfinie dans son ensemble, à l'échelle globale. Il me semble qu'il faut pour cela creuver l'abcès, et en débattre ouvertement, tout autant qu'échanger sur le rôle des chercheu.r.se.s face aux crises sociales et environnementales (voir Box~\hyperref[box:Conferences]{\uppercase\expandafter{A\relax}.\theConferences{}}).



\begin{spacing}{1.5}
~
\end{spacing}

\boxspenoref{\thecounterboxspe{}}{Frais de publication}{
\setlength\parindent{12pt}
Les frais de publication de mon premier papier de thèse m'ont paru faramineux, plus de 5800 euros pour un article publié en accès libre, dans un journal \textit{Biogéosciences}, supposé bon marché. Cela m'a révolté, et m'a incité lors de la soumission de mon deuxième papier de thèse à \textit{Earth System Dynamics} à demander une réduction des frais de publication avec un petit message au regard du caractère inacceptable à mes yeux de tirer profit économiquement des publications scientifiques :

\begin{spacing}{0.7}
~
\end{spacing}

\centerline{\begin{minipage}{0.85\textwidth}
\setlength\parindent{12pt}
\textit{I am aware that APCs will be levied for the publication of my paper, but I would like to apply for a discount or waiver.I am applying for a discount of 30~\%. Justification: I do not think it is acceptable to make a profit from scientific productions. It would therefore seem fair and honest to me, in view of the values of science and the service it is supposed to render to society, to pay the cost price of my article, as well as the small margin (10~\% according to your documentation) assigned to authors lacking the financial sources to pay for the article processing charges. So I am asking for a 30~\% discount, estimating your margin at 30~\%. If you are honest enough to share your actual margin percentage, then I will be open to a request for a smaller reduction up to your actual margin - which I hope is less than 30~\% all the same. In the hope that you will understand my request, I would like to thank you or making your journals open access. I do think that this request should not be a request but a form of normalcy and responsibility of research in front of society.}
\end{minipage}}

\begin{spacing}{0.7}
~
\end{spacing}

Cela n'a évidemment rien donné, je m'y attendais, et je n'ai pas poursuivi, car ma thèse était sur le point de se terminer. Mais une fois le papier publié, je l'espère, je n'hésiterai pas à les recontacter, peut-être d'une manière différente, mais je n'hésiterai pas, à nouveau, à déranger ce qui me semble intolérable.

\label{box:Publication}
}
\\

\begin{spacing}{1.5}
~
\end{spacing}

%%%%%%%%%%%%%%%%%%%%%%%%%%%%%%%%%%%%%%%%%%%%%%
{\noindent\Large\textbf{Se (dé)battre}}
\\

Pourquoi vouloir que la nouvelle génération de chercheu.r.se.s reproduise les erreurs du passé ? Et en sciences du climat, comment peut-on produire autant de connaissances sans même leur donner vie, de la valeur en adaptant notre manière de faire dans le contexte professionnel ? Comment accepter notre contradiction ? Car, bien que désagréable, ce que notre recherche nous apprend est aussi une invitation aujourd'hui, une contrainte malgré nous peut-être demain, à repenser notre activité de recherche à la hauteur de l'urgence climatique, entre autres. Mais c'est une marche conséquente à gravir, ensemble, car projeter le changement climatique ne nous engage en rien, au contraire de l'action pour y faire face. S'engager résulte de la réponse à un "Pourquoi ?", à l'intégration d'informations qui nous mettent en mouvement. S'engager c'est prendre des risques pour soi, mais investir son énergie pour le bien commun. Néanmoins face au système de la recherche actuel, s'engager peut isoler, rejeter. Et pourtant quand je commence à soulever la question du "Comment ?", j'en reviens toujours au besoin de partager, d'échanger pour faire le deuil d'un passé indésirable et se mettre en marche pour un futur désirable. Mais face à une question d'ordre systémique, se battre, s'apparente plutôt à se débattre à première vue, on est pris par l'immensité des changements à entreprendre. Alors j'essaie d'accepter de toujours faire trop peu face au défi qui s'offre à nous, tenter au moins, participer à l'effort. J'essaie notamment de trouver la forme d'engagement qui me convient, qui me fait avancer, en tentant de définir les échelles spatiales, temporelles et sociales de mon combat. 
\\


%%%%%%%%%%%%%%%%%%%%%%%%%%%%%%%%%%%%%%%%%%%%%%
{\noindent\Large\textbf{Se rabattre ou rebattre les cartes}}
\\


Chaque matin, le même réveil, celui qui sonne pour me mettre debout et enfourcher mon vélo, y retourner. En route pour une routine qui tourne à vide, car chaque jour, la même interrogation, celle qui questionne notre fonctionnement. Alors, chaque soir, confronté à un sens-unique sur le trajet du chemin le plus court pour rentrer, la question d'emprunter un autre chemin au carrefour qui le précède se pose, mais sans jamais donner suite. Chaque nuit, ça me travaille, j'y repense. Mais le lendemain, happé par une charge de travail, le sens trouvé pendant la nuit est piétiné, et je pars à nouveau à contre-sens en reprenant le même sens unique que la veille, l'avant-veille. Une fois encore, c'est ignorer ce que ce carrefour interroge, le système de la recherche tel qu'il est, et nos habitudes. Il symbolise les tourments que je peux avoir face à une recherche dont le modèle repose sur des valeurs et des principes qui ne me correspondent pas, même si l'activité en elle-même, est passion. Lorsque je réalise que pour transformer son modèle, il faudrait changer fondamentalement notre modèle de société, je suis saisis d'une certaine peur. Mais renoncer à quelque chose, c'est aussi aspirer à autre chose. De la peur à l'audace, pourquoi se rabattre, emprunter à nouveau le même sens unique et abandonner l'idée de transformer le système de la recherche, plutôt que de rebattre les cartes, faire un détour et chercher à le redéfinir ? À quoi bon continuer ainsi quand on peut changer de voix ?
\\


%%%%%%%%%%%%%%%%%%%%%%%%%%%%%%%%%%%%%%%%%%%%%%
{\noindent\Large\textbf{Ouverture}}
\\

La recherche, lorsqu'elle ne répond pas à une simple curiosité, est un service pour la société, un investissement pour l'avenir. Perdurer une recherche productiviste, consumériste, fournissant la matière au système tel que nous le connaissons aujourd'hui, sans en questionner en amont son fonctionnement, son orientation, ses productions et leurs impacts, est un non-sens. À la folie d'apprendre par l'expérience, une sagesse réflexive finira par répondre, j'y crois. Je rêve d'un système de la recherche, notamment en sciences du climat, qui montre l'exemple. S'engager pour repenser la recherche est un premier pas vers l'action. Alors, je m'obstine non pas à refuser mais à apprendre, non pas à rejeter mais à accepter, non pas à renoncer mais à aspirer. Et si l'on me demande si je souhaite continuer dans la recherche, je répondrais que j'aime chercher, que je continuerai de chercher avec passion, et que, si je pouvais faire de ma recherche une oeuvre d'art, alors je le ferais.
\\


























%%%%%%%%%%%%%%%%%%%%%%%%%%%%%%%%%%%%%%%%%%%%%%
